% Emacs, this is -*-latex-*-

\ex{Conceitos de grafos direcionados}
\label{ex:directed-graph-concepts}
Considere o seguinte grafo direcionado:

\begin{center}
    \scalebox{1}{ % x,y
        \begin{tikzpicture}[dgraph]
            \node[cont] (a) at (0,2) {$a$};
            \node[cont] (z) at (2,2) {$z$};
            \node[cont] (q) at (1,1) {$q$};
            \node[cont] (e) at (1,-0.4) {$e$};
            \node[cont] (h) at (3,1) {$h$};
            \draw(a) -- (q);
            \draw(z) -- (h);
            \draw(z) -- (q);
            \draw(q) -- (e);
    \end{tikzpicture}}
\end{center}

\begin{exenumerate}
    \item Liste todas as trilhas no grafo (de comprimento m\'aximo)
    
    \begin{solution}
        Temos
        $$(a,q,e) \quad \quad (a,q,z,h) \quad \quad  (h,z,q,e)$$
        e as correspondentes com os n\'os de in\'icio e fim trocados.
    \end{solution}
    
    \item Liste todos os caminhos direcionados no grafo (de comprimento m\'aximo)
    \begin{solution}
        $(a,q,e) \quad \quad (z,q,e) \quad \quad (z,h) $
    \end{solution}
    
    \item Quais s\~ao os descendentes de $z$?
    
    \begin{solution}
        $\desc(z) =\{q,e,h\}$
    \end{solution}
    
    
    \item Quais s\~ao os n\~ao-descendentes de $q$?
    
    \begin{solution}
        $\nondesc(q) = \{a,z,h,e\} \setminus \{e\} = \{a,z,h\}$
    \end{solution}
    
    
    \item Quais das seguintes ordena\'c\~oes s\~ao topol\'ogicas para o grafo?
    \begin{itemize}
        \item (a,z,h,q,e)
        \item (a,z,e,h,q)
        \item (z,a,q,h,e)
        \item (z,q,e,a,h)
    \end{itemize}
    
    \begin{solution}
        \begin{itemize}
            \item (a,z,h,q,e): sim
            \item (a,z,e,h,q): n\~ao ($q$ \'e pai de $e$ e, portanto, deve vir antes de $e$ na ordena\'c\~ao)
            \item (z,a,q,h,e): sim
            \item (z,q,e,a,h): n\~ao ($a$ \'e pai de $q$ e, portanto, deve vir antes de $q$ na ordena\'c\~ao)
        \end{itemize}
    \end{solution}
    
\end{exenumerate}

% --------------------------------------------------- 

\ex{Conex\~oes can\^onicas}
\label{ex:canonical-connections}
Aqui derivamos as independ\^encias que ocorrem nas tr\^es conex\~oes can\^onicas que existem em DAGs, mostradas na Figura \ref{fig:canonical-connections}.

\begin{figure}[h!]
    \begin{center}
        \subfloat[Conex\~ao serial]{
            \scalebox{0.9}{ % x,y
                \begin{tikzpicture}[dgraph]
                    \node[cont] (x) at (0,0) {$x$};
                    \node[cont] (z) at (2,0) {$z$};
                    \node[cont] (y) at (4,0) {$y$};
                    \draw(x) -- (z);
                    \draw(z) -- (y);
            \end{tikzpicture}}
        }\hspace{2ex}
        \subfloat[Conex\~ao divergente]{
            \scalebox{0.9}{ % x,y
                \begin{tikzpicture}[dgraph]
                    \node[cont] (x) at (0,0) {$x$};
                    \node[cont] (z) at (2,0) {$z$};
                    \node[cont] (y) at (4,0) {$y$};
                    \draw(z) -- (x);
                    \draw(z) -- (y);
            \end{tikzpicture}}
        }\hspace{2ex}
        \subfloat[Conex\~ao convergente]{
            \scalebox{0.9}{ % x,y
                \begin{tikzpicture}[dgraph]
                    \node[cont] (x) at (0,0) {$x$};
                    \node[cont] (z) at (2,0) {$z$};
                    \node[cont] (y) at (4,0) {$y$};
                    \draw(x) -- (z);
                    \draw(y) -- (z);
            \end{tikzpicture}}
        }
    \end{center}
    \caption{\label{fig:canonical-connections}As tr\^es conex\~oes can\^onicas em DAGs.}
\end{figure}

\begin{exenumerate}
    \item Para a conex\~ao serial, use a propriedade de Markov ordenada para mostrar que 
    $x \independent y \,|\, z$.
    \begin{solution}
        A \'unica ordena\'c\~ao topol\'ogica \'e $x, z, y$. Os predecessores de
        $y$ s\~ao $\pre_y = \{x, z\}$ e seus pais $\pa_y = \{z\}$. A
        propriedade de Markov ordenada
        \begin{equation}
            y \independent ( \pre_y \setminus \pa_y ) \mid \pa_y
        \end{equation}
        torna-se, assim, $y \independent (\{x, z\} \setminus z) \mid z$. Portanto, temos
        \begin{equation}
            y \independent x \mid z,
        \end{equation}
        o que \'e o mesmo que $x \independent y \mid z$, j\'a que a
        rela\'c\~ao de independ\^encia \'e sim\'etrica.
        
        Isso significa que se o estado ou valor de $z$ for conhecido (ou seja, se
        a vari\'avel aleat\'oria $z$ for ``instanciada''), a evid\^encia sobre $x$
        n\~ao mudar\'a nossa cren\'ca sobre $y$, e vice-versa. Dizemos que
        o n\'o $z$ est\'a ``fechado'' e que a trilha entre $x$ e $y$
        est\'a ``bloqueada'' pelo $z$ instanciado. Em outras palavras, conhecer
        o valor de $z$ bloqueia o fluxo de evid\^encia \emph{entre} $x$
        e $y$.
        
    \end{solution}
    \item Para a conex\~ao serial, mostre que a marginal $p(x,y)$ geralmente n\~ao se fatoriza em $p(x)p(y)$, ou seja, que $x \independent y$ n\~ao se sustenta.
    \begin{solution}
        Existem v\'arias maneiras de mostrar o resultado. Uma delas \'e
        apresentar um exemplo onde a independ\^encia n\~ao
        ocorre. Considere, por exemplo, o seguinte modelo:
        \begin{align}
            x & \sim \normal(x; 0, 1)\label{eq:serial-x}\\
            z &= x + n_z\\
            y &= z + n_y \label{eq:serial-y}
        \end{align}
        onde $n_z \sim \normal(n_z; 0, 1)$ e $n_y \sim \normal(n_y;
        0, 1)$, ambos sendo estatisticamente independentes de $x$. Aqui
        $\normal(\cdot; 0, 1)$ denota a fdp Gaussiana com m\'edia 0 e
        vari\^ancia 1, e $x \sim \normal(x; 0, 1)$ significa que amostramos
        $x$ da distribui\'c\~ao $\normal(x; 0, 1)$. Assim, $p(z|x) =
        \normal(z; x, 1)$, $p(y|z) = \normal(y; z, 1)$ e $p(x, y,z) =
        p(x) p(z|x) p(y|z) = \normal(x; 0, 1) \normal(z; x,
        1)\normal(y; z, 1)$.
        
        Embora pud\'essemos manipular as fdps para mostrar o resultado, \'e
        aqui mais f\'acil trabalhar com o modelo generativo nas Equa\'c\~oes
        \eqref{eq:serial-x} a \eqref{eq:serial-y}. Eliminando $z$
        das equa\'c\~oes, ao substituir a defini\'c\~ao de $z$ em
        \eqref{eq:serial-y}, temos
        \begin{equation}
            y = x + n_z + n_y,
        \end{equation}
        que descreve a distribui\'c\~ao marginal de $(x, y)$. Vemos
        que $\E[xy]$ \'e
        \begin{align}
            \E[xy] & = \E[x^2 + xn_z + x n_y]\\
            & =  \E[x^2] + \E[x]\E[n_z] + \E[x]\E[n_y]\\
            & = 1 + 0 + 0
        \end{align}
        onde usamos a linearidade da esperan\'ca, que $x$ \'e
        independente de $n_z$ e $n_y$, e que $x$ tem m\'edia
        zero. Se $x$ e $y$ fossem independentes (ou apenas n\~ao correlacionados),
        ter\'iamos $\E[xy] = \E[x]\E[y] = 0$. No entanto, como $\E[xy] \neq \E[x]
        \E[y]$, $x$ e $y$ n\~ao s\~ao independentes.
        
        Em termos simples, isso significa que se o estado de $z$ for
        desconhecido, ent\~ao a evid\^encia ou informa\'c\~ao sobre $x$ influenciar\'a
        nossa cren\'ca sobre $y$, e vice-versa. A evid\^encia pode
        fluir atrav\'es de $z$ entre $x$ e $y$. Dizemos que o n\'o $z$
        est\'a ``aberto'' e a trilha entre $x$ e $y$ est\'a ``ativa''.
        
    \end{solution}
    \item Para a conex\~ao divergente, use a propriedade de Markov ordenada
    para mostrar que $x \independent y \,|\, z$.
    \begin{solution}
        Uma ordena\'c\~ao topol\'ogica \'e $z, x, y$. Os predecessores de
        $y$ s\~ao $\pre_y = \{x, z\}$ e seus pais $\pa_y = \{z\}$. A
        propriedade de Markov ordenada
        \begin{equation}
            y \independent ( \pre_y \setminus \pa_y ) \mid \pa_y
        \end{equation}
        torna-se, novamente,
        \begin{equation}
            y \independent x \mid z,
        \end{equation}
        que \'e, dado que a rela\'c\~ao de independ\^encia \'e sim\'etrica, o
        mesmo que $x \independent y \mid z$.
        
        Como na conex\~ao serial, se o estado ou valor $z$ for conhecido,
        a evid\^encia sobre $x$ n\~ao mudar\'a nossa cren\'ca sobre $y$, e vice-
        versa. Conhecer $z$ fecha o n\'o $z$, o que bloqueia a trilha
        entre $x$ e $y$.
    \end{solution}
    
    \item Para a conex\~ao divergente, mostre que a marginal $p(x,y)$
    geralmente n\~ao se fatoriza em $p(x)p(y)$, ou seja, que $x \independent y$ n\~ao se
    sustenta.
    
    \begin{solution}
        Assim como para a conex\~ao serial, basta dar um exemplo
        onde $x \independent y$ n\~ao ocorra. Consideramos o seguinte
        modelo generativo:
        \begin{align}
            z & \sim \normal(z; 0, 1)\\
            x &= z + n_x\\
            y &= z + n_y
        \end{align}
        onde $n_x \sim \normal(n_x; 0, 1)$ e $n_y \sim \normal(n_y; 0,
        1)$, e eles s\~ao independentes entre si e das outras
        vari\'aveis. Temos $\E[x] = \E[z + n_x] = \E[z]+\E[n_x]=0$. Por
        outro lado,
        \begin{align}
            \E[x y] & = \E[(z + n_x)( z + n_y)]\\
            & = \E[z^2 + z(n_x+n_y) + n_xn_y]\\
            & = \E[z^2] +  \E[z(n_x+n_y)] +  \E[n_xn_y]\\
            & = 1 + 0 + 0
        \end{align}
        Portanto, $\E[xy] \neq \E[x]\E[y]$ e n\~ao temos que $x
        \independent y$ se sustente.
        
        Em uma conex\~ao divergente, como na conex\~ao serial, se
        o estado de $z$ for desconhecido, ent\~ao a evid\^encia ou informa\'c\~ao sobre
        $x$ influenciar\'a nossa cren\'ca sobre $y$, e vice-
        versa. A evid\^encia pode fluir atrav\'es de $z$ entre $x$ e $y$. Dizemos
        que o n\'o $z$ est\'a aberto e a trilha entre $x$ e $y$ est\'a
        ativa.
    \end{solution}
    
    \item Para a conex\~ao convergente, mostre que $x \independent y$.
    
    \begin{solution}
        Podemos aqui usar novamente a propriedade de Markov ordenada com a
        ordena\'c\~ao $y, x, z$. Como $\pre_x = \{y\}$ e $\pa_x =
        \varnothing$, temos
        \begin{equation}
            x \independent ( \pre_x \setminus \pa_x ) \mid \pa_x = x \independent y.
        \end{equation}
        Alternativamente, podemos usar a defini\'c\~ao b\'asica de modelos
        gr\'aficos direcionados, i.e.,
        \begin{align}
            p(x,y,z) & = k(x) k(y) k(z \mid x, y)
        \end{align}
        juntamente com o resultado de que os kernels (fatores) s\~ao fdps/fmps
        (condicionais) v\'alidas e iguais \`as condicionais/marginais
        em rela\'c\~ao \`a distribui\'c\~ao conjunta $p(x,y,z)$, i.e.,
        \begin{align}
            k(x) & = p(x)\\
            k(y) &= p(y) \\
            k(z | x, y) & = p(z|x,y) \quad \text{\small (n\~ao necess\'ario na prova abaixo)}
        \end{align}
        Integrar em rela\'c\~ao a $z$ resulta em
        \begin{align}
            p(x,y) & = \int p(x,y,z) \ud z\\
            & =  \int  k(x) k(y) k(z \mid x, y) \ud z\\
            & = k(x) k(y) \underbrace{\int k(z \mid x, y) \ud z}_{1}\\
            & = p(x) p(y)
        \end{align}
        Portanto, $p(x,y)$ se fatoriza em suas marginais, o que significa que $x
        \independent y$.
        
        Assim, quando n\~ao temos evid\^encia sobre $z$, a evid\^encia sobre $x$
        n\~ao mudar\'a nossa cren\'ca sobre $y$, e vice-versa. Para a
        conex\~ao convergente, se nenhuma evid\^encia sobre $z$ estiver dispon\'ivel, o
        n\'o $z$ est\'a fechado, o que bloqueia a trilha entre $x$ e $y$.
        
    \end{solution}
    
    \item Para a conex\~ao convergente, mostre que $x \independent y \mid
    z$ geralmente n\~ao se sustenta.
    
    \begin{solution}
        Damos um exemplo simples onde $x \independent y \mid
        z$ n\~ao ocorre.
        
        Considere
        \begin{align}
            x & \sim \normal(x;0, 1)\\ y & \sim \normal(y;0, 1)\\ z&= xy +
            n_z
        \end{align}
        onde $n_z \sim \normal(n_z; 0, 1)$, independente das outras
        vari\'aveis. Da \'ultima equa\'c\~ao, temos
        \begin{equation}
            x y = z - n_z
        \end{equation}
        Temos, portanto,
        \begin{align}
            \E[xy \mid z] & = \E[z - n_z \mid z] \\ &= z - 0
        \end{align}
        Por outro lado, $\E[x y] = \E[x]\E[y] = 0$. Como $\E[xy \mid z] \neq \E[x y]$, $x \independent y \mid z$ n\~ao pode ocorrer.
        
        A intui\'c\~ao aqui \'e que se voc\^e conhece o valor do produto
        $xy$, mesmo que sujeito a ru\'ido, conhecer o valor de $x$ permite
        que voc\^e estime o valor de $y$ e vice-versa.
        
        Mais geralmente, para conex\~oes convergentes, se a evid\^encia ou
        informa\'c\~ao sobre $z$ estiver dispon\'ivel, a evid\^encia sobre $x$
        influenciar\'a a cren\'ca sobre $y$, e vice-versa. Dizemos que
        a informa\'c\~ao sobre $z$ abre o n\'o $z$, e a evid\^encia pode fluir
        entre $x$ e $y$.
        
        Nota: informa\'c\~ao sobre $z$ significa que \emph{$z$ ou um de seus
            descendentes} \'e observado, veja o exerc\'icio \ref{ex:independencies2}.
        
    \end{solution}
    
\end{exenumerate}

% --------------------------------------------------- 


\ex{Propriedades de Markov ordenada e local, d-separa\'c\~ao}
\label{ex:ordered-local-d-sep-independencies01}
Continuamos com a investiga\'c\~ao do grafo do \exref{ex:directed-graph-concepts} mostrado abaixo para refer\^encia.

\begin{center}
    \scalebox{1}{ % x,y
        \begin{tikzpicture}[dgraph]
            \node[cont] (a) at (0,2) {$a$};
            \node[cont] (z) at (2,2) {$z$};
            \node[cont] (q) at (1,1) {$q$};
            \node[cont] (e) at (1,-0.4) {$e$};
            \node[cont] (h) at (3,1) {$h$};
            \draw(a) -- (q);
            \draw(z) -- (h);
            \draw(z) -- (q);
            \draw(q) -- (e);
    \end{tikzpicture}}
\end{center}

\begin{exenumerate}
    \item A ordena\'c\~ao $(z,h,a,q,e)$ \'e topol\'ogica para o grafo. Quais s\~ao as independ\^encias que seguem da propriedade de Markov ordenada?
    
    \begin{solution}
        Uma distribui\'c\~ao que se fatoriza sobre o grafo satisfaz as independ\^encias
        $$ x_i \independent \left(\pre_i \setminus \pa_i\right) \mid \pa_i \text{ para todo } i$$
        para todas as ordena\'c\~oes das vari\'aveis que s\~ao topol\'ogicas para o grafo. A ordena\'c\~ao entra em jogo atrav\'es dos conjuntos de predecessores $\pre_i = \{x_1, \ldots, x_{i-1}\}$ das vari\'aveis $x_i$; o grafo atrav\'es dos conjuntos de pais $\pa_i$.
        
        Para o grafo e a ordena\'c\~ao topol\'ogica especificada, os conjuntos de predecessores s\~ao
        $$\pre_z = \varnothing, \pre_h = \{z\}, \pre_a = \{z,h\}, \pre_q = \{z,h,a\}, \pre_e =\{z,h,a,q\}$$
        Os conjuntos de pais dependem apenas do grafo e n\~ao da ordena\'c\~ao topol\'ogica. Eles s\~ao:
        $$\pa_z = \varnothing, \pa_h = \{z\}, \pa_a = \varnothing, \pa_q =\{a,z\}, \pa_e = \{q\}, $$
        A propriedade de Markov ordenada l\^e-se como $x_i \independent \left(\pre_i \setminus \pa_i \right)\mid \pa_i$ onde os $x_i$ referem-se \`as vari\'aveis ordenadas, por exemplo, $x_1 = z, x_2 = h, x_3 = a,$ etc.
        
        Com
        $$ \pre_h \setminus \pa_h = \varnothing \quad \pre_a \setminus \pa_a = \{z,h\} \quad \pre_q \setminus \pa_q = \{h\} \quad \pre_e \setminus \pa_e = \{z,h,a\}$$
        obtemos, assim,
        $$ h \independent \varnothing \mid z \quad \quad
        a \independent \{z,h\} \quad \quad q \independent
        h \mid \{a,z\} \quad \quad e \independent \{z,h,a\} \mid q$$ A
        rela\'c\~ao $ h \independent \varnothing \mid z$ deve ser entendida como
        ``n\~ao h\'a vari\'avel da qual $h$ seja independente dado $z$'' e
        deve, portanto, ser removida da lista. Note que podemos possivelmente
        obter mais rela\'c\~oes de independ\^encia para vari\'aveis que ocorrem mais tarde na
        ordena\'c\~ao topol\'ogica. Isso ocorre porque o conjunto
        $\pre \setminus \pa$ s\'o pode aumentar quando o conjunto de predecessores
        $\pre$ torna-se maior.
    \end{solution}
    
    \item Quais s\~ao as independ\^encias que seguem da propriedade de Markov local?
    
    \begin{solution}
        Os n\~ao-descendentes s\~ao
        $$\nondesc(a) = \{z,h\} \quad \nondesc(z) = \{a\} \quad \nondesc(h) = \{a,z,q,e\}$$
        $$\quad \nondesc(q) = \{a,z,h\} \quad \nondesc(e) = \{a,q,z,h\}$$
        Com os conjuntos de pais como anteriormente, as independ\^encias que seguem da propriedade de Markov local s\~ao $x_i \independent \left( \nondesc(x_i)\setminus \pa_i \right) \mid \pa_i$, ou seja:
        $$a \independent \{z,h\} \quad \quad z \independent a \quad \quad h \independent \{a,q,e\} \mid z \quad \quad q \independent  h \mid \{a,z\} \quad \quad e \independent \{a,z,h\} \mid q$$
        
        
    \end{solution}
    
    
    \item As rela\'c\~oes de independ\^encia obtidas via propriedade de Markov ordenada e local
    incluem $q \independent h \mid \{a,z\}$. Verifique
    a independ\^encia usando d-separa\'c\~ao.
    
    \begin{solution}
        
        A \'unica trilha de $q$ para $h$ passa por $z$, que est\'a em uma configura\'c\~ao cauda-cauda. Como $z$ faz parte do conjunto de condicionamento, a trilha est\'a bloqueada e o resultado segue.
    \end{solution}
    
    \item Use d-separa\'c\~ao para verificar se $a \independent h \mid e$ se sustenta.
    \begin{solution}
        A trilha de $a$ para $h$ \'e mostrada abaixo em vermelho, juntamente com os estados padr\~ao dos n\'os ao longo da trilha.
        \begin{center}
            \scalebox{1}{ % x,y
                \begin{tikzpicture}[dgraph]
                    \node[cont] (a) at (0,2) {$a$};
                    \node[cont] (z) at (2,2) {$z$};
                    \node[cont] (q) at (1,1) {$q$};
                    \node[cont] (e) at (1,-0.4) {$e$};
                    \node[cont] (h) at (3,1) {$h$};
                    \draw (a) -- (q);
                    \draw[color=red,-,style=dashed] (a) -- (q);
                    \draw(z) -- (h);
                    \draw[color=red,-,style=dashed] (z) -- (h);
                    \draw(z) -- (q);
                    \draw[color=red,-,style=dashed](z) -- (q);
                    \draw(q) -- (e);
                    \node[left=0.01of q] {fechado};
                    \node[above=0.01of z] {aberto};
            \end{tikzpicture}}
        \end{center}
        Condicionar em $e$ abre o n\'o $q$, j\'a que $q$ est\'a em uma configura\'c\~ao de colisor no caminho. 
        \begin{center}
            \scalebox{1}{ % x,y
                \begin{tikzpicture}[dgraph]
                    \node[cont] (a) at (0,2) {$a$};
                    \node[cont] (z) at (2,2) {$z$};
                    \node[cont] (q) at (1,1) {$q$};
                    \node[contobs] (e) at (1,-0.4) {$e$};
                    \node[cont] (h) at (3,1) {$h$};
                    \draw (a) -- (q);
                    \draw[color=red,-,style=dashed] (a) -- (q);
                    \draw(z) -- (h);
                    \draw[color=red,-,style=dashed] (z) -- (h);
                    \draw(z) -- (q);
                    \draw[color=red,-,style=dashed](z) -- (q);
                    \draw(q) -- (e);
                    \node[left=0.01of q] {aberto};
                    \node[above=0.01of z] {aberto};
            \end{tikzpicture}}
        \end{center}
        A trilha de $a$ para $h$ est\'a, portanto, ativa, o que significa que a rela\'c\~ao n\~ao se sustenta porque $a \notind h \mid e$ para algumas distribui\'c\~oes que se fatorizam sobre o grafo.
    \end{solution}
    
    
    \item Assuma que todas as vari\'aveis no grafo sejam bin\'arias. Quantos n\'umeros voc\^e precisa especificar, ou aprender com os dados, para especificar totalmente a distribui\'c\~ao de probabilidade?
    \begin{solution}
        O grafo define um conjunto de fun\'c\~oes de massa de probabilidade (fmp) que se fatorizam como
        $$ p(a,z,q,h,e) = p(a) p(z) p(q|a,z) p(h|z) p(e|q)$$
        Para especificar um membro do conjunto, precisamos especificar as fmps (condicionais) do lado direito. As fmps (condicionais) podem ser vistas como tabelas, e o n\'umero de elementos que precisamos especificar nas tabelas s\~ao:\\
        - 1 para $p(a)$ \\
        - 1 para $p(z)$ \\
        - 4 para $p(q | a,z)$ \\
        - 2 para $p(h | z)$\\
        - 2 para $p(e|q)$ \\
        No total, h\'a 10 n\'umeros para especificar. Isso contrasta com $2^5-1 = 31$ para uma distribui\'c\~ao sem independ\^encias.
        Note que o n\'umero de par\^ametros a serem especificados poderia ser ainda mais reduzido atrav\'es de hip\'oteses param\'etricas. 
        
    \end{solution}
    
\end{exenumerate}


\ex{Mais sobre as propriedades de Markov ordenada e local, d-separa\'c\~ao}
Continuamos com a investiga\'c\~ao do grafo abaixo:

\begin{center}
    \scalebox{1}{ % x,y
        \begin{tikzpicture}[dgraph]
            \node[cont] (a) at (0,2) {$a$};
            \node[cont] (z) at (2,2) {$z$};
            \node[cont] (q) at (1,1) {$q$};
            \node[cont] (e) at (1,-0.4) {$e$};
            \node[cont] (h) at (3,1) {$h$};
            \draw(a) -- (q);
            \draw(z) -- (h);
            \draw(z) -- (q);
            \draw(q) -- (e);
    \end{tikzpicture}}
\end{center}

\begin{exenumerate}
    \item Por que a propriedade de Markov ordenada ou local n\~ao pode ser usada para verificar se $a \independent h \mid e$ pode ocorrer? 
    \begin{solution}
        As independ\^encias que seguem da propriedade de Markov ordenada ou local exigem o condicionamento em conjuntos de pais. No entanto, $e$ n\~ao \'e pai de nenhum n\'o, de modo que a asser\'c\~ao de independ\^encia acima n\~ao pode ser verificada via propriedade de Markov ordenada ou local. 
    \end{solution}
    
    \item As rela\'c\~oes de independ\^encia obtidas via propriedade de Markov ordenada e local incluem $a \independent \{z,h\}$. Verifique a independ\^encia usando d-separa\'c\~ao.
    
    \begin{solution}
        Todos os caminhos de $a$ para $z$ ou $h$ passam pelo n\'o $q$, que forma uma conex\~ao cabe\'ca-cabe\'ca ao longo dessa trilha. Como nem $q$ nem seu descendente $e$ fazem parte do conjunto de condicionamento, a trilha est\'a bloqueada e a rela\'c\~ao de independ\^encia segue.
    \end{solution}
    
    
    \item Determine o Markov blanket (cobertura de Markov) de $z$.
    \begin{solution}
        O Markov blanket \'e dado pelos pais, filhos e copais (outros pais dos filhos). Assim: $\MB(z) = \{a,q,h\}$.
    \end{solution}
    
    
    \item Verifique que $ q \independent  h \mid \{a,z\}$ se sustenta ao manipular a distribui\'c\~ao de probabilidade induzida pelo grafo.
    \begin{solution}
        
        Uma defini\'c\~ao b\'asica de independ\^encia estat\'istica condicional $x_1
        \independent x_2 \mid x_3$ \'e que a conjunta (condicional)
        $p(x_1, x_2 \mid x_3)$ seja igual ao produto das marginais
        (condicionais) $p(x_1 \mid x_3)$ e $p(x_2 \mid x_3)$. Em outras
        palavras, para vari\'aveis aleat\'orias discretas,
        \begin{align}
            x_1 &\independent x_2 \mid x_3 &  &\Longleftrightarrow &  p(x_1, x_2 \mid x_3) &= \left( \sum_{x_2}  p(x_1, x_2 \mid x_3) \right) \left( \sum_{x_1}  p(x_1, x_2 \mid x_3) \right)
        \end{align}
        Assim, respondemos \`a quest\~ao mostrando que (use integrais em
        caso de vari\'aveis aleat\'orias cont\'inuas)
        \begin{align}
            p(q,h | a,z) &= \left(\sum_h  p(q,h | a,z)\right) \left( \sum_q  p(q,h | a,z)\right)
        \end{align}
        
        Primeiro, note que o grafo define um conjunto de fun\'c\~oes de densidade
        ou massa de probabilidade que se fatorizam como
        $$ p(a,z,q,h,e) = p(a) p(z) p(q|a,z) p(h|z) p(e|q)$$ Usamos ent\~ao a regra da soma para
        computar a distribui\'c\~ao conjunta de $(a,z,q,h)$, ou seja, a
        distribui\'c\~ao de todas as vari\'aveis que ocorrem em $p(q,h|a,z)$ 
        \begin{align}
            p(a,z,q,h) &= \sum_{e} p(a,z,q,h,e)\\
            &= \sum_{e}  p(a) p(z) p(q|a,z) p(h|z) p(e|q)\\
            & =  p(a) p(z) p(q|a,z) p(h|z) \underbrace{\sum_{e}  p(e|q)}_{1}\\
            & =  p(a) p(z) p(q|a,z) p(h|z),
        \end{align}
        onde $\sum_{e}  p(e|q) = 1$ porque as fdps/fmps (condicionais) s\~ao normalizadas de modo que integram/somam para um. Temos ainda
        \begin{align}
            p(a,z) & = \sum_{q,h} p(a,z,q,h)\\
            & = \sum_{q,h} p(a) p(z) p(q|a,z) p(h|z)\\
            & = p(a) p(z) \sum_{q}  p(q|a,z) \sum_h p(h|z)\\
            & = p(a) p(z)
        \end{align}
        de modo que
        \begin{align}
            p(q,h | a,z) & = \frac{ p(a,z,q,h)}{ p(a,z)} \\
            & = \frac{ p(a) p(z) p(q|a,z) p(h|z)}{ p(a) p(z)}\\
            & =  p(q|a,z) p(h|z).
        \end{align}        
        Vemos ainda que $p(q| a,z)$ e $p(h|z)$ s\~ao as marginais de $p(q,h|a,z)$, ou seja,
        \begin{align}
            p(q| a,z) & = \sum_h  p(q,h | a,z)\\
            p(h|z) & = \sum_q  p(q,h | a,z).
        \end{align}
        Isso significa que
        \begin{equation}
            p(q,h | a,z) =\left( \sum_h  p(q,h | a,z) \right) \left(\sum_q  p(q,h | a,z) \right),
        \end{equation}
        o que mostra que $q \independent h | a,z$.
        
        Vemos que usar o grafo para determinar a independ\^encia \'e
        mais f\'acil do que manipular a fmp/fdp.
        
    \end{solution}
    
\end{exenumerate}

% --------------------------------------------------------------- 


\ex{Cl\'inica de t\'orax {\small \citep[baseado em][Exerc\'icio 3.3]{Barber2012}}} 
\label{ex:chest-clinic}
O modelo gr\'afico direcionado na Figura \ref{fig:chest-clinic} \'e sobre
o diagn\'ostico de doen\'ca pulmonar (t=tuberculose ou l=c\^ancer de pulm\~ao). Neste
modelo, acredita-se que uma visita a algum lugar ``$a$'' aumente a
probabilidade de tuberculose.

\begin{figure}[h!]
    \centering
    \includegraphics[width=0.9\textwidth]{chest-clinic}
    \caption{\label{fig:chest-clinic}Modelo gr\'afico para o \exref{ex:chest-clinic} (Barber Figura 3.15).}
\end{figure}

\begin{exenumerate}
    \item Explique qual das seguintes rela\'c\~oes de independ\^encia se sustenta para todas as distribui\'c\~ao que se fatorizam sobre o grafo.
    \begin{enumerate}
        \item $t \independent s \mid d$
        
        \begin{solution}
            \begin{itemize}
                \item Existem duas trilhas de $t$ para $s$: $(t,e,l,s)$ e $(t,e,d,b,s)$. 
                \item A trilha $(t,e,l,s)$ apresenta um n\'o colisor $e$ que \'e aberto pela vari\'avel de condicionamento $d$. A trilha est\'a, portanto, ativa e n\~ao precisamos verificar a segunda trilha, pois para a independ\^encia todas as trilhas precisariam estar bloqueadas.
                \item A rela\'c\~ao de independ\^encia, portanto, geralmente n\~ao se sustenta.
            \end{itemize}
            
        \end{solution}
        
        \item $l \independent b \mid s$
        
        \begin{solution}
            \begin{itemize}
                \item Existem duas trilhas de $l$ para $b$: $(l,s,b)$ e $(l,e,d,b)$
                \item A trilha $(l,s,b)$ est\'a bloqueada por $s$ ($s$ est\'a em uma configura\'c\~ao cauda-cauda e faz parte do conjunto de condicionamento)
                \item A trilha $(l,e,d,b)$ est\'a bloqueada pela configura\'c\~ao de colisor para o n\'o $d$.
                \item Todas as trilhas est\~ao bloqueadas, de modo que a rela\'c\~ao de independ\^encia se sustenta.
            \end{itemize}
        \end{solution}
        
    \end{enumerate}
    
    \item Podemos simplificar $p(l|b,s)$ para $p(l|s)$?
    
    \begin{solution}
        Como $l \independent b \mid s$, temos $p(l|b,s) = p(l|s)$.
    \end{solution}
    
\end{exenumerate}


\ex{Mais sobre a cl\'inica de t\'orax {\small \citep[baseado em][Exerc\'icio 3.3]{Barber2012}}} 
\label{ex:chest-clinic2}
Considere o modelo gr\'afico direcionado na Figura \ref{fig:chest-clinic}.

\begin{exenumerate}
    \item Explique qual das seguintes rela\'c\~oes de independ\^encia se sustenta para todas as distribui\'c\~ao que se fatorizam sobre o grafo.
    \begin{enumerate}
        \item $a \independent s \mid l$
        \begin{solution}
            \begin{itemize}
                \item Existem duas trilhas de $a$ para $s$: $(a,t,e,l,s)$ e $(a,t,e,d,b,s)$
                \item A trilha $(a,t,e,l,s)$ apresenta um n\'o colisor $e$ que bloqueia a trilha (a trilha tamb\'em \'e bloqueada por $l$).
                \item A trilha $(a,t,e,d,b,s)$ \'e bloqueada pelo n\'o colisor $d$. 
                \item Todas as trilhas est\~ao bloqueadas, de modo que a rela\'c\~ao de independ\^encia se sustenta.
            \end{itemize}
        \end{solution}
        
        \item $a \independent s \mid l, d$
        
        \begin{solution}
            \begin{itemize}
                \item Existem duas trilhas de $a$ para $s$: $(a,t,e,l,s)$ e $(a,t,e,d,b,s)$
                \item A trilha $(a,t,e,l,s)$ apresenta um n\'o colisor $e$ que \'e aberto pela vari\'avel de condicionamento $d$, mas o n\'o $l$ \'e fechado pela vari\'avel de condicionamento $l$: a trilha est\'a bloqueada.
                \item A trilha $(a,t,e,d,b,s)$ apresenta um n\'o colisor $d$ que \'e aberto pelo condicionamento em $d$. Nesta trilha, $e$ n\~ao est\'a em uma configura\'c\~ao cabe\'ca-cabe\'ca (colisor), de modo que todos os n\'os est\~ao abertos e a trilha ativa.
                \item Assim, a rela\'c\~ao de independ\^encia geralmente n\~ao se sustenta.
            \end{itemize}
        \end{solution}           
        
    \end{enumerate}
    
    \item Seja $g$ uma fun\'c\~ao (determin\'istica) de $x$ e $t$. O valor esperado $\E[g(x,t) \mid l,b]$ \'e igual a $\E[g(x,t) \mid l]$?
    
    \begin{solution}
        A quest\~ao resume-se a verificar se $x,t \independent b \mid l$. Para que a rela\'c\~ao de independ\^encia ocorra, todas as trilhas tanto de $x$ quanto de $t$ para $b$ precisam estar bloqueadas por $l$.
        
        \begin{itemize}
            \item Para $x$, temos as trilhas $(x,e,l,s,b)$ e $(x,e,d,b)$
            \item A trilha $(x,e,l,s,b)$ \'e bloqueada por $l$.
            \item A trilha $(x,e,d,b)$ \'e bloqueada pela configura\'c\~ao de colisor do n\'o $d$.
            \item Para $t$, temos as trilhas $(t,e,l,s,b)$ e $(t,e,d,b)$
            \item A trilha $(t,e,l,s,b)$ \'e bloqueada por $l$.
            \item A trilha $(t,e,d,b)$ \'e bloqueada pela configura\'c\~ao de colisor do n\'o $d$.
        \end{itemize}
        Como todas as trilhas est\~ao bloqueadas, temos $x,t \independent b \mid l$ e $\E[ g(x,t) \mid l,b] = \E[ g(x,t) \mid l]$.
        
    \end{solution}
\end{exenumerate}


\ex{Modelos ocultos de Markov (HMM)}
\label{ex:HMM}
Este exerc\'icio \'e sobre modelos gr\'aficos direcionados que s\~ao especificados pelo seguinte DAG:
\begin{center}
    \scalebox{0.8}{
        \begin{tikzpicture}[dgraph]
            \node[cont] (y1) at (0,0) {$y_1$};
            \node[cont] (y2) at (2,0) {$y_2$};
            \node[cont] (y3) at (4,0) {$y_3$};
            \node[cont] (y4) at (6,0) {$y_4$};
            \node[cont] (x1) at (0,2) {$x_1$};
            \node[cont] (x2) at (2,2) {$x_2$};
            \node[cont] (x3) at (4,2) {$x_3$};
            \node[cont] (x4) at (6,2) {$x_4$};
            \draw(x1)--(y1);\draw(x2)--(y2);\draw(x3)--(y3);\draw(x4)--(y4);
            \draw(x1)--(x2);\draw(x2)--(x3);\draw(x3)--(x4);
    \end{tikzpicture}}
\end{center}
Esses modelos s\~ao chamados de modelos de Markov ``ocultos'' porque normalmente assumimos que apenas observamos os $y_i$ e n\~ao os $x_i$, que seguem um modelo de Markov.

\begin{exenumerate}
    \item Mostre que todos os modelos probabil\'isticos especificados pelo DAG se fatorizam como
    $$p(x_1,y_1, x_2, y_2, \ldots,x_4,y_4) = p(x_1)p(y_1|x_1)p(x_2 | x_1)p(y_2|x_2)p(x_3|x_2)p(y_3|x_3)p(x_4|x_3)p(y_4|x_4)$$
    \begin{solution}
        Da defini\'c\~ao de modelos gr\'aficos direcionados segue que
        $$p(x_1,y_1, x_2, y_2, \ldots,x_4,y_4) = \prod_{i=1}^4 p(x_i |
        \pa(x_i)) \prod_{i=1}^4 p(y_i | \pa(y_i)).$$ O resultado \'e ent\~ao
        obtido ao notar que o pai de $y_i$ \'e dado por $x_i$ para todo $i$,
        e que o pai de $x_i$ \'e $x_{i-1}$ para $i=2,3,4$ e que $x_1$ n\~ao possui pai ($\pa(x_1) = \varnothing$).
    \end{solution}
    %
    \item Derive as independ\^encias implicadas pela propriedade de Markov ordenada com a ordena\'c\~ao topol\'ogica $(x_1, y_1, x_2, y_2, x_3, y_3, x_4, y_4)$
    
    \begin{solution}
        $$y_i \independent x_1, y_1, \ldots, x_{i-1}, y_{i-1} \mid x_i \quad \quad x_i \independent x_1, y_1, \ldots, x_{i-2},y_{i-2},y_{i-1} \mid x_{i-1}$$
        
    \end{solution}
    
    \item Derive as independ\^encias implicadas pela propriedade de Markov ordenada com a ordena\'c\~ao topol\'ogica $(x_1, x_2, \ldots, x_4, y_1, \ldots, y_4)$.
    \begin{solution} Para os $x_i$, usamos que para $i \ge 2$: $\pre(x_i) = \{x_1, \ldots, x_{i-1}\}$ e $\pa(x_i) = x_{i-1}$. Para os $y_i$, usamos que
        $\pre(y_1) = \{x_1,\ldots, x_4\}$, que $\pre(y_i) = \{x_1,\ldots, x_4, y_1, \ldots, y_{i-1}\}$ for $i>1$, e que $\pa(y_i) = x_i$. A propriedade de Markov ordenada ent\~ao resulta em:
        \begin{align*}
            x_3 \independent x_1 \mid x_2 && x_4 \independent \{x_1,x_2\} \mid x_3\\
            y_1 \independent \{x_2, x_3, x_4\} \mid x_1 &&  y_2 \independent \{x_1, x_3, x_4, y_1\} \mid x_2 \\
            y_3 \independent \{x_1, x_2, x_4, y_1, y_2\} \mid x_3 &&  y_4 \independent \{x_1, x_2, x_3, y_1, y_2, y_3\} \mid x_4
        \end{align*}
        
    \end{solution}
    
    \item $y_4 \independent y_1 \mid y_3$ se sustenta?
    
    \begin{solution}
        A trilha $y_1-x_1-x_2-x_3-x_4-y_4$ est\'a ativa: nenhum dos n\'os
        est\'a em uma configura\'c\~ao de colisor, de modo que seu estado padr\~ao \'e
        aberto e condicionar em $y_3$ n\~ao bloqueia nenhum dos n\'os na
        trilha.
        
        Embora $x_1-x_2-x_3-x_4$ forme uma cadeia de Markov, onde por exemplo $x_4
        \independent x_1 \mid x_3$ se sustenta, isso n\~ao ocorre para a distribui\'c\~ao
        dos $y$'s.
        
    \end{solution}
\end{exenumerate}


\ex{Caracteriza\'c\~ao alternativa de independ\^encias}
\label{ex:independencies1}
Vimos que $x \independent y | z$ \'e caracterizado por
$p(x,y | z) =p(x | z) p(y| z)$
ou, equivalentemente, por
$p(x| y, z) = p(x | z)$.
Mostre que outras caracteriza\'c\~oes equivalentes s\~ao
\begin{align}
    p(x,y,z) &= p(x|z) p(y|z) p(z) \quad \text{e} \label{eq:characterisation1-alt} \\
    p(x,y,z) &= a(x,z) b(y,z) \quad \text{para algumas fun\'c\~oes n\~ao-neg. \;} a(x,z) \text{\; e \;} b(x,z). \label{eq:characterisation2-alt} 
\end{align}
A caracteriza\'c\~ao na Equa\'c\~ao \eqref{eq:characterisation2-alt} \'e particularmente importante para modelos gr\'aficos n\~ao direcionados.
\begin{solution}
    Primeiro mostramos a equival\^encia de $p(x,y | z) =p(x | z) p(y| z)$ e
    $p(x,y,z) = p(x|z) p(y|z) p(z)$: Pela regra do produto, temos
    $$p(x,y,z) = p(x,y|z) p(z).$$ Se $p(x,y | z) =p(x | z) p(y| z)$,
    segue que $p(x,y,z) = p(x|z) p(y|z) p(z)$. Para mostrar a dire\'c\~ao oposta, assuma que $p(x,y,z) = p(x|z) p(y|z) p(z)$ se sustenta. Por
    compara\'c\~ao com a decomposi\'c\~ao na regra do produto, segue-se
    que devemos ter $p(x,y | z) = p(x | z) p(y| z)$ sempre que $p(z)>0$
    (basta considerar este caso porque para $z$ onde $p(z)=0$, $p(x,y | z)$ pode n\~ao estar
    univocamente definido em primeiro lugar).
    
    A Equa\'c\~ao \eqref{eq:characterisation1-alt} implica
    \eqref{eq:characterisation2-alt} com $a(x,z) = p(x | z)$ e $b(y,z) =
    p(y| z)p(z)$. Agora mostramos a inversa. Vamos assumir que $p(x,y,z)
    = a(x,z) b(y,z)$. Pela regra do produto, temos
    \begin{align}
        p(x,y | z)p(z) & = a(x,z) b(y,z).\\
    \end{align}
    Somar sobre $y$ resulta em
    \begin{align}
        \sum_y  p(x,y | z)p(z) & = p(z) \sum_y p(x,y | z)\\
        & = p(z) p(x|z)
    \end{align}
    Al\'em disso,
    \begin{align}
        \sum_y  p(x,y | z)p(z) & = \sum_y  a(x,z) b(y,z)\\
        & = a(x,z) \sum_y b(y,z)
    \end{align}
    de modo que 
    \begin{equation}
        a(x,z) = \frac{p(z) p(x | z)}{\sum_y b(y,z)}
        \label{eq:a-alt}
    \end{equation}
    Como a soma de $p(x | z)$ sobre $x$ \'e igual a um, temos
    \begin{equation}
        \sum_x a(x,z) = \frac{p(z)}{\sum_y b(y,z)}.
        \label{eq:asum-alt}
    \end{equation}
    
    Agora, somar $p(x,y | z) p(z)$ sobre $x$ fornece
    \begin{align}
        \sum_x  p(x,y | z) p(z) & =  p(z) \sum_x p(x,y|z). \\
        & = p(y | z) p(z)
    \end{align}
    Tamb\'em temos
    \begin{align}
        \sum_x  p(x,y | z) p(z) &= \sum_x a(x,z) b(y,z) \\
        &= b(y,z) \sum_x a(x,z) \\
        &\stackrel{\eqref{eq:asum-alt}}{=} b(y,z) \frac{p(z)}{\sum_y b(y,z)}
    \end{align}
    de modo que
    \begin{equation}
        p(y | z) p(z) =  p(z) \frac{b(y,z)}{\sum_y b(y,z)}
        \label{eq:brel-alt}
    \end{equation}
    Temos, portanto,
    \begin{align}
        p(x,y,z) & = a(x,z) b(y,z) \\
        &  \stackrel{\eqref{eq:a-alt}}{=} \frac{p(z) p(x | z)}{\sum_y b(y,z)} b(y,z) \\
        & = p(x | z)p(z) \frac{ b(y,z)}{\sum_y b(y,z)} \\
        & \stackrel{\eqref{eq:brel-alt}}{=} p(x|z) p(y|z) p(z)
    \end{align}
    que \'e a Equa\'c\~ao \eqref{eq:characterisation1-alt}.
\end{solution}


\ex{Mais sobre independ\^encias}
\label{ex:independencies2-bis}
Este exerc\'icio \'e sobre propriedades e caracteriza\'c\~oes adicionais de independ\^encia estat\'istica.

\begin{exenumerate}
    \item Sem usar d-separa\'c\~ao, mostre que $x \independent \{y,w\} \mid z$ implica que $x \independent y \mid z$ e $x \independent w \mid z$.\\
    \emph{Dica: use a defini\'c\~ao de independ\^encia estat\'istica em termos da fatoriza\'c\~ao de fmps/fdps.}
    
    \begin{loesung}
        Consideramos a distribui\'c\~ao conjunta $p(x,y,w | z)$. Por hip\'otese,
        \begin{equation}
            p(x,y,w | z) = p(x |z) p(y,w | z)
        \end{equation}
        Temos que mostrar que $x \independent y| z$ e $x \independent w | z$. Por simplicidade, assumimos que as vari\'aveis t\^em valores discretos. Se n\~ao, substitua a soma abaixo por uma integral.
        
        Para mostrar que $x \independent y| z$, marginalizamos $p(x,y,w | z)$
        sobre $w$ para obter
        \begin{align}
            p(x,y | z ) & = \sum_w p(x,y,w | z) \\
            & = \sum_w  p(x |z) p(y,w | z) \\
            & = p(x | z) \sum_w p(y,w | z)
        \end{align}
        Como $\sum_w p(y,w | z)$ \'e a marginal $p(y |z)$, temos
        \begin{equation}
            p(x,y | z ) =  p(x | z) p(y | z),
        \end{equation}
        o que significa que $x \independent y | z$.
        
        Para mostrar que $x \independent w| z$, marginalizamos de forma similar
        $p(x,y,w | z)$ sobre $y$ para obter $p(x,w | z) =
        p(x|z) p(w |z)$, o que significa que $x \independent w | z$.
        
    \end{loesung}
    
    
    \item Para o modelo gr\'afico direcionado abaixo, mostre que as seguintes duas afirma\'c\~oes se sustentam sem usar d-separa\'c\~ao:
    
    \begin{align}
        &x \independent y \quad \text{e} \label{eq:statement1-alt} \\
        &x \notind y \mid w \label{eq:statement2-alt}
    \end{align}
    
    \begin{center}
        \scalebox{0.9}{ % x,y
            \begin{tikzpicture}[dgraph]
                \node[cont] (x) at (0,0) {$x$};
                \node[cont] (z) at (2,0) {$z$};
                \node[cont] (w) at (2,-1.5) {$w$};
                \node[cont] (y) at (4,0) {$y$};
                \draw(x) -- (z);
                \draw(y) -- (z);
                \draw(z) -- (w);
        \end{tikzpicture}}
    \end{center}
    O exerc\'icio mostra que n\~ao apenas condicionar em um n\'o colisor, mas tamb\'em em um de seus descendentes, ativa a trilha entre $x$ e $y$. Voc\^e pode usar o resultado de que $x \independent y |w \Leftrightarrow p(x,y,w) = a(x,w)b(y,w)$ para algumas fun\'c\~oes n\~ao-negativas $a(x,w)$ e $b(y,w)$.
    
    \begin{solution}
        O modelo gr\'afico corresponde \`a fatoriza\'c\~ao $$p(x,y,z,w) = p(x) p(y) p(z|x,y) p(w|z).$$
        Para a marginal $p(x,y)$, temos que somar (integrar) sobre todos os $(z,w)$:
        \begin{align}
            p(x,y) &= \sum_{z,w} p(x,y,z,w) \\
            & = \sum_{z,w} p(x) p(y) p(z|x,y) p(w|z)\\
            & = p(x) p(y) \sum_{z,w} p(z| x,y) p(w|z)\\
            & = p(x) p(y) \underbrace{\sum_{z} p(z| x,y)}_{1} \underbrace{\sum_{w} p(w|z)}_{1}\\
            & = p(x) p(y)
        \end{align}
        Como $p(x,y) = p(x) p(y)$, temos $x \independent y$.
        
        Para $x \notind y | w$, compute $p(x,y,w)$ e use o resultado $x \independent y |w \Leftrightarrow p(x,y,w) = a(x,w)b(y,w)$.
        \begin{align}
            p(x,y,w) & = \sum_z p(x,y,z,w)\\
            & = \sum_z p(x) p(y) p(z|x,y) p(w|z) \\
            & = p(x) \underbrace{p(y) \sum_z p(z|x,y) p(w|z)}_{k(x,y,w)}
        \end{align}
        Como $p(x,y,w)$ n\~ao pode ser fatorado como $a(x,w) b(y,w)$, a rela\'c\~ao $x \independent y | w$ n\~ao pode ocorrer de forma geral.
    \end{solution}
    
\end{exenumerate}


\ex{Independ\^encias em modelos gr\'aficos direcionados}

Considere o seguinte grafo ac\'iclico direcionado:
\begin{center}
    \scalebox{0.9}{ % x,y
        \begin{tikzpicture}[dgraph]
            \node[cont] (x1) at (0,0) {$x_1$};
            \node[cont] (x2) at (0,-2) {$x_2$};
            \node[cont] (x3) at (0,-4) {$x_3$};
            \node[cont] (x4) at (2,-1.5) {$x_4$};
            \node[cont] (x5) at (2,-3) {$x_5$};
            \node[cont] (x6) at (2,-4.5) {$x_6$};
            \node[cont] (x7) at (4,0) {$x_7$};
            \node[cont] (x8) at (4,-2) {$x_8$};
            \node[cont] (x9) at (6,-1) {$x_9$};
            
            \draw (x1) -- (x2);
            \draw (x2) -- (x3);
            \draw (x1) -- (x4);
            \draw (x4) -- (x5);
            \draw (x5) -- (x6);
            \draw (x7) -- (x4);
            \draw (x7) -- (x8);
            \draw (x3) -- (x6);
            \draw (x8) -- (x6);
            \draw (x7) -- (x9);
    \end{tikzpicture}}
\end{center}
Para cada uma das afirma\'c\~oes abaixo, determine se ela se sustenta para
todos os modelos probabil\'isticos que se fatorizam sobre o grafo. Forne\'ca uma
justificativa para sua resposta.

\begin{exenumerate}
    \item $p(x_7 | x_2) = p(x_7)$
    
    \begin{solution}
        Sim, se sustenta. $x_2$ \'e um n\~ao-descendente de $x_7$,
        $\pa(x_7)=\varnothing$ e, portanto, pela propriedade de Markov local, $x_7$
        $\independent x_2$, de modo que $p(x_7 | x_2) = p(x_7)$.
    \end{solution}
    
    \item $x_1 \notind x_3$  
    
    \begin{solution}
        N\~ao, n\~ao se sustenta. $x_1$ e $x_3$ s\~ao d-conectados, o que apenas implica independ\^encia para \emph{algumas} e n\~ao todas as distribui\'c\~oes que se fatorizam sobre o grafo. O grafo geralmente s\'o nos permite ler independ\^encias e n\~ao depend\^encias.
        
    \end{solution}
    
    
    \item $p(x_1,x_2,x_4) \propto \phi_1(x_1,x_2) \phi_2(x_1,x_4)$ para algumas fun\'c\~oes n\~ao-negativas $\phi_1$ e $\phi_2$. 
    
    \begin{solution}
        Sim, se sustenta. A afirma\'c\~ao \'e equivalente a $x_2 \independent x_4 \mid x_1$. Existem tr\^es trilhas de $x_2$ para $x_4$, que est\~ao todas bloqueadas:
        \begin{enumerate}
            \item $x_2-x_1-x_4$: esta trilha est\'a bloqueada porque $x_1$ est\'a em uma conex\~ao cauda-cauda e \'e observado, o que fecha o n\'o.
            \item $x_2-x_3-x_6-x_5-x_4$: esta trilha est\'a bloqueada porque
            $x_3, x_6, x_5$ est\~ao em uma configura\'c\~ao de colisor, e $x_6$
            n\~ao \'e observado (e n\~ao possui nenhum descendente).
            \item $x_2-x_3-x_6- x_8- x_7- x_4$: esta trilha est\'a bloqueada porque
            $x_3, x_6, x_8$ est\~ao em uma configura\'c\~ao de colisor, e $x_6$
            n\~ao \'e observado (e n\~ao possui nenhum descendente).
        \end{enumerate}
        Portanto, pela propriedade de Markov global (d-separa\'c\~ao), a independ\^encia se sustenta.
    \end{solution}
    
    
    \item $x_2 \independent x_9 \mid \{x_6, x_8\}$ 
    
    \begin{solution}
        N\~ao, n\~ao se sustenta. Condicionar em $x_6$ abre o n\'o colisor $x_4$ na trilha $x_2 - x_1 - x_4 - x_7 - x_9$, de modo que a trilha fica ativa.
    \end{solution}
    
    
    \item $x_8 \independent \{x_2,x_9\} \mid \{x_3,x_5, x_6, x_7\}$ 
    
    \begin{solution}
        Sim, se sustenta. $\{x_3,x_5,x_6,x_7\}$ \'e o Markov blanket de $x_8$, de modo que $x_8$ \'e independente dos n\'os restantes dado o Markov blanket.
    \end{solution}
    
    
    \item $\E[ x_2 \cdot x_3 \cdot x_4 \cdot x_5 \cdot x_8 \mid x_7] = 0$ se $\E[x_8 \mid x_7] = 0$
    
    \begin{solution}
        Sim, se sustenta. $\{x_2,x_3,x_4,x_5\}$ s\~ao n\~ao-descendentes de $x_8$, e $x_7$ \'e o pai de $x_8$, de modo que $x_8 \independent \{x_2,x_3,x_4,x_5\} \mid x_7$. Isso significa que $\E[ x_2 \cdot x_3 \cdot x_4 \cdot x_5 \cdot x_8 \mid x_7] = \E[ x_2 \cdot x_3 \cdot x_4 \cdot x_5 \mid x_7] \E[x_8 \mid x_7] =0$.
    \end{solution}
    
    
\end{exenumerate}


\ex{Independ\^encias em modelos gr\'aficos direcionados}
Considere o seguinte grafo ac\'iclico direcionado:
\begin{center}
    \begin{tikzpicture}[dgraph]
        
        \node[cont] (m1) at (0,0) {$m_1$};
        \node[cont] (s1) at (2,0) {$s_1$};
        \node[cont] (u1) at (0,-1.5) {$u_1$};
        \node[cont] (v1) at (2,-1.5) {$v_1$};
        \node[cont] (x1) at (0,-3) {$x_1$};
        \node[cont] (y1) at (2,-3) {$y_1$};
        \node[cont] (theta1) at (1,-4) {$\theta_1$};
        
        \draw (m1) -- (u1);
        \draw (m1) -- (v1);
        \draw (s1) -- (u1);
        \draw (s1) -- (v1);
        \draw (u1) -- (x1);
        \draw (v1) -- (y1);
        \draw (theta1) -- (x1);
        \draw (theta1) -- (y1);
        
        
        \node[cont] (m2) at (4,0) {$m_2$};
        \node[cont] (s2) at (6,0) {$s_2$};
        \node[cont] (u2) at (4,-1.5) {$u_2$};
        \node[cont] (v2) at (6,-1.5) {$v_2$};
        \node[cont] (x2) at (4,-3) {$x_2$};
        \node[cont] (y2) at (6,-3) {$y_2$};
        \node[cont] (theta2) at (5,-4) {$\theta_2$};
        
        \draw (m2) -- (u2);
        \draw (m2) -- (v2);
        \draw (s2) -- (u2);
        \draw (s2) -- (v2);
        \draw (u2) -- (x2);
        \draw (v2) -- (y2);
        \draw (theta2) -- (x2);
        \draw (theta2) -- (y2);
        
        \draw (theta1) -- (theta2);
        
    \end{tikzpicture}
\end{center}

Para cada uma das afirma\'c\~ao abaixo, determine se ela se sustenta para
todos os modelos probabil\'isticos que se fatorizam sobre o grafo. Forne\'ca uma
justificativa para sua resposta.

\begin{exenumerate}
    \item $x_1 \independent x_2$
    
    \begin{solution}
        N\~ao se sustenta. A trilha $x_1 - \theta_1 - \theta_2 - x_2$ est\'a
        ativa (n\~ao bloqueada) porque nenhum dos n\'os est\'a em uma configura\'c\~ao
        de colisor ou no conjunto de condicionamento.
    \end{solution}
    
    \item $p(x_1, y_1, \theta_1, u_1) \propto \phi_A(x_1, \theta_1, u_1)\phi_B(y_1, \theta_1, u_1)$ para algumas fun\'c\~oes n\~ao-negativas $\phi_A$ e $\phi_B$ 
    
    \begin{solution}
        Se sustenta. A afirma\'c\~ao \'e equivalente a $x_1 \independent y_1$
        $\mid \{\theta_1, u_1\}$. O conjunto de condicionamento $\{\theta_1,
        u_1\}$ bloqueia todas as trilhas de $x_1$ para $y_1$ porque ambos
        est\~ao apenas em configura\'c\~oes seriais em todas as trilhas de $x_1$ para
        $y_1$; logo, a independ\^encia se sustenta pela propriedade de Markov
        global. Justificativa alternativa: o conjunto de condicionamento \'e
        o Markov blanket de $x_1$, e $x_1$ e $y_1$ n\~ao s\~ao
        vizinhos, o que implica a independ\^encia.
    \end{solution}
    
    \item $v_2 \independent \{u_1, v_1, u_2, x_2\} \mid \{m_2, s_2, y_2, \theta_2\}$
    
    \begin{solution}
        Se sustenta. O conjunto de condicionamento \'e o Markov blanket de $v_2$
        (o conjunto de pais, filhos e copais): o conjunto de
        pais \'e $\pa(v_2)=\{m_2, s_2\}$, $y_2$ \'e o \'unico filho de
        $v_2$, e $\theta_2$ \'e o \'unico outro pai de
        $y_2$. E $v_2$ \'e independente de todas as outras vari\'aveis
        dado o seu Markov blanket.
    \end{solution}
    
    \item $\E [ m_2 \mid  m_1 ] = \E[m_2]$
    
    \begin{solution}
        Se sustenta. Existem quatro trilhas de $m_1$ para $m_2$, a saber, via
        $x_1$, via $y_1$, via $x_2$, via $y_2$. Em todas as trilhas as quatro
        vari\'aveis est\~ao em uma configura\'c\~ao de colisor, de modo que cada uma das
        trilhas est\'a bloqueada. Pela propriedade de Markov global
        (d-separa\'c\~ao), isso significa que $m_1 \independent m_2$, o que
        implica que $\E[m_2 \mid m_1] = \E[m_2]$.
        
        Justificativa alternativa 1: $m_2$ \'e um n\~ao-descendente de $m_1$
        e $\pa(m_2) = \varnothing$. Pela propriedade de Markov local direcionada,
        uma vari\'avel \'e independente de seus n\~ao-descendentes
        dados os pais, logo $m_2 \independent m_1$.
        
        Justificativa alternativa 2: Podemos escolher uma ordena\'c\~ao
        topol\'ogica onde $m_1$ e $m_2$ s\~ao as duas primeiras
        vari\'aveis. Al\'em disso, seus conjuntos de pais est\~ao vazios. Pela
        propriedade de Markov ordenada direcionada, temos, portanto, $m_1 \independent m_2$.
        
    \end{solution}
    
    
\end{exenumerate}